Why does a bounded sequence of functions sometimes refuse to converge, and what is left when it does? In infinite dimension three comfortable rules of linear algebra fail, and the repairs built between 1900 and 1935 are among the most useful ideas in mathematics.

This book proves them. Hahn--Banach, the three theorems that follow from Baire's theorem, and weak compactness. Hilbert spaces. The spectral theorem for matrices, and the one for compact self-adjoint operators. Graph spectra, expanders, and the Riemann hypothesis for graphs. Fourier analysis on the line, on distributions, on finite groups and their fast algorithms, and on locally compact groups, with Haar measure and the $p$-adic numbers. Then the complex analysis this series needs: Cauchy's theorem with the winding number, the growth of entire functions, and the Mellin transform. Tauberian theorems read a sequence off its transform and prove the prime number theorem again. A last chapter of problems turns those tools on the values of the zeta function.

The spectral theorem proved here is the one for matrices and for compact self-adjoint operators. Six appendices hold the linear algebra behind the main path. Every exercise has a hint and a solution that shows how the argument is found.
